% vim: ai sw=2 spell spelllang=cs

\section{Komunikace jádro $\leftrightarrow$ uživatelský prostor}
\frame{\sectionpage}

\begin{frame}
  \frametitle{Komunikace I.}
  \begin{enumerate}
    \item Voláním funkce: system call (syscall)

    \pause

    \begin{itemize}
      \item Skok do jádra speciální instrukcí (x86\_64: \code{syscall})
      \begin{itemize}
	\item O skok se stará libc (\code{fwrite}$\rightarrow$\code{write}$\rightarrow$\code{syscall}$\rightarrow$instrukce)
	\item Drahá operace (přepnutí kontextu)
      \end{itemize}

      \pause

      \item Do registru se uloží číslo operace
      \begin{itemize}
	\item V libc: převodní ,,tabulka'' funkce$\rightarrow$číslo
	\item V jádře: převodní tabulka číslo$\rightarrow$funkce
      \end{itemize}

      \pause

      \item \code{syscall(__NR_fork)}
      \item Demo: \code{__NR_fork} v \file{unistd.h}; příklady z \file{pb173/02/syscalls/}
    \end{itemize}

    \pause

    \item Speciální syscall
    \begin{itemize}
      \item Knihovna přilinkovaná jádrem
      \item \file{vdso.so}
    \end{itemize}
  \end{enumerate}

  \pause
  \medskip

  \begin{center}
  \heading{Každá nová funkcionalita = nový syscall}
  \end{center}

\end{frame}

\begin{frame}{Úkol}

  \heading{V userspace: pomocí \textit{funkce} \code{syscall} vypište nějakou
    informaci}
  \begin{enumerate}
    \item Inspirujte se v \file{pb173/02/syscalls/c.c}
    \item Použijte \code{syscall} a \code{__NR_write}
  \end{enumerate}

  \medskip

  Je třeba znát prototypy funkcí.
  \begin{itemize}
    \item Většinou jsou dokumentované: \doc{man write}
    \item Jinak použít lxr a najít implementaci v jádře
  \end{itemize}

\end{frame}

\begin{frame}
  \frametitle{Komunikace II.}

  \begin{enumerate}
    \setcounter{enumi}{2}

    \item Speciální soubory v \file{/dev}
    \begin{itemize}
      \item Komunikace přes soubor
      \begin{itemize}
	\item Standardní operace -- \code{open}, \code{read}, \dots
	\item Tj.\ není nutný nový syscall
      \end{itemize}

      \pause

      \item Identifikované trojicí blokové/znakové, major a minor číslo
      \begin{itemize}
	\item Většinou major=ovladač, minor=zařízení (tty: c, maj 4, min 0--63)
      \end{itemize}

      \pause

      \item Blokové (disky apod.)
      \begin{itemize}
	\item Komunikace po blocích
	\item Nebudeme se jimi zabývat (popsány v LDD)
      \end{itemize}

      \pause

      \item Znakové (ostatní)
      \begin{itemize}
	\item Komunikace po znacích (bajtech)
	\item \emph{Zbytek \iffimuni tohoto cvičení \else této sekce \fi}
      \end{itemize}

      \pause

      \item Seznam rezervovaných v \doc{Documentation/devices.txt} a
	\file{/proc/devices}
    \end{itemize}

    \pause

    \item Sockety, roury, \dots
  \end{enumerate}

\end{frame}

\mysection{Znaková zařízení}

\begin{frame}[fragile]
  \frametitle{Obsluhované funkce}

  \code{struct file_operations} (\header{linux/fs.h}, LDD3 3. kapitola)
  \begin{itemize}
  \item Funkce jako:
  \begin{lstlisting}
ssize_t (*write)(struct file *file, const char __user *buf, size_t count, loff_t *offp)
  \end{lstlisting}

  \pause

  \item \code{file->private\_data} slouží programátorovi (libovolně)
  \item \code{offp} slouží programátorovi (k poznamenání průběhu)
  \item \code{\_\_user} značí ukazatel od uživatele (\emph{POZOR} tomu nevěříme)

  \pause

  \item Návratové hodnoty (podle návratového typu funkce)
  \begin{itemize}
    \item \code{int} -- záporné = \texttt{-Echyba}, jinak 0
    \item \code{ssize\_t} -- záporné = \texttt{-Echyba}, jinak počet
      zpracovaných znaků
    \item Vždy podle standardu (např.\ POSIX)
  \end{itemize}
  \end{itemize}
\end{frame}

\begin{frame}{Úkol}

  \heading{Práce s lxr}
  \begin{enumerate}
  \item Najděte všechny možné funkce, které lze obsluhovat
  \begin{itemize}
    \item Najděte \code{struct file_operations} v lxr
    \item Projděte strukturu
  \end{itemize}

  \item Najděte možné chybové návratové hodnoty
  \begin{itemize}
    \item Najděte \code{EPERM} v lxr
    \item Projděte ostatní
    \item Podívejte se do \doc{man write} na možné návratové hodnoty
  \end{itemize}
  \end{enumerate}
\end{frame}

\begin{frame}[fragile]{Příklad operací}

  \begin{lstlisting}
struct file_operations my_fops = {
  .owner = THIS_MODULE,
  .open = my_open,
  .write = my_write,
};
  \end{lstlisting}
  \pause
  \begin{lstlisting}
int my_open(struct inode *inode, struct file *filp) {
  filp->private_data = 3;
  return 0;
}
  \end{lstlisting}
  \pause
  \begin{lstlisting}
ssize_t my_write(struct file *filp, const char __user *buf, size_t count, loff_t *offp) {
  /* filp->private_data is 3 here */
  if (count == 0)
    return -EINVAL;
  *offp += count; /* remember position for the next time */
  return count; /* we "wrote" all bytes */
}
  \end{lstlisting}

\end{frame}

\begin{frame}
  \frametitle{Úkol}

  \heading{Obsluha \code{open}, \code{read}, \code{write}, \code{release} (tj.
    close)}

  \begin{enumerate}
    \item Definice \code{struct file_operations}
    \item Vytvořte funkce dle prototypu z \code{file_operations} (lxr)
    \begin{itemize}
      \item Open a release s nějakým \code{printk} a \code{return}
      \item Read a write prozatím jen s \code{return}
      \item Návratové hodnoty
      \begin{itemize}
	\item Open a release vracejí 0 (žádná chyba)
	\item Read též (0 znamená EOF)
	\item Write vrací \code{count} (zapsáno vše)
      \end{itemize}

    \end{itemize}

    \item Jen zkuste přeložit
    \begin{itemize}
      \item Nemá smysl vkládat
    \end{itemize}
  \end{enumerate}

\end{frame}

\begin{frame}
  \frametitle{Vytvoření znakového zařízení}

  \begin{itemize}
  \item LDD3 3.\ a 6.\ kapitola
  \item 2--3 kroky

  \begin{enumerate}
  \item Registrace rozsahu major+minor (v \code{module_init})
    \begin{itemize}
    \item Buď \code{alloc_chrdev_region} (chci čísla přidělit), anebo
      \code{register_chrdev_region} (čísla mám), nakonec
      \code{unregister_chrdev_region} (\header{linux/fs.h})
    \item Přidání záznamu do \file{/proc/devices}
    \end{itemize}

  \pause

  \item Registrace jednotlivých minorů (při připojení zařízení)
    \begin{itemize}
    \item \code{cdev_init}, \code{cdev_add}, \code{cdev_del}
      (\header{linux/cdev.h})
    \item Parametr pro \code{cdev_init} je \code{struct file_operations}
    \item Po odpovídajícím \cmd{mknod} lze zařízení používat
    \end{itemize}

  \pause

  \item Podat zprávu \cmd{udev} (vytvoření \file{/dev/*}) -- nepovinné
    \begin{itemize}
    \item \code{device_create}, \code{device_destroy}
      (\header{linux/device.h})
    \item Předem je potřeba vytvořit \code{class} (v \code{module_init})
    \end{itemize}
  \end{enumerate}

  \end{itemize}

\end{frame}

\begin{frame}{Úkol}

  \heading{Vytvoření znakového zařízení}
  \begin{enumerate}
  \item Definujte si globální \code{struct cdev} (znakové zařízení)

  \item V \code{module_init}
  \begin{itemize}
    \item Zavolejte \code{alloc_chrdev_region(&dev, 0, 10, "jmeno")}
    \begin{itemize}
      \item \code{dev} je návratová hodnota typu \code{dev_t}
      \item Vypište si \code{MAJOR(dev)} a \code{MINOR(dev)}
    \end{itemize}
    \item Inicializujte svou globální \code{struct cdev} (\code{cdev_init})
    \item Přidejte svoje \code{struct cdev} do systému (\code{cdev_add} a
      \code{dev} shora)
  \end{itemize}

  \item V \code{module_exit}
  \begin{itemize}
    \item \code{cdev_del}
    \item \code{unregister_chrdev_region}
  \end{itemize}

  \item \cmd{make} a \cmd{insmod}
  \item Ověřte vytvoření v \file{/proc/devices}
  \item Podle výpisu z 2.\ bodu: \cmd{mknod soubor c MAJ MIN}
  \item Vyzkoušejte \cmd{cat soubor}

  \end{enumerate}
\end{frame}

\begin{frame}
  \frametitle{Uživatelská paměť}

  \begin{itemize}
    \item Něco, čemu nelze věřit (NULL, ukazatel do tabulek oprávnění, \dots)
    \item \emph{POZOR}: nutnost kontroly

    \pause

    \item \code{copy_from_user}, \code{copy_to_user}
    \begin{itemize}
      \item ,,\code{memcpy}'' s kontrolou
      \item Vracejí počet \emph{NEzkopírovaných} znaků (tj.\ 0 znamená OK)
    \end{itemize}

    \pause

    \item \code{get_user}, \code{put_user}
    \begin{itemize}
      \item ,,\code{var = *buf}'' a ,,\code{*buf = var}'' s kontrolou
      \item Tzn.\ jen pro primitiva (\code{char}, \code{short}, \code{int},
	\code{long})
      \item Vracejí 0 nebo chybu (záporná hodnota)
    \end{itemize}

    \item Definované v \header{linux/uaccess.h}
  \end{itemize}

  \pause

  \heading{Demo:} \file{pb173/02/}

\end{frame}

\begin{frame}
  \frametitle{Úkol}

  \heading{Dopsat těla funkcí \code{read} a \code{write} tak, aby zpracovávala
    data}

  \begin{enumerate}

  \item \code{write} vypíše uživatelský buffer pomocí \code{printk}
  \begin{itemize}
    \item \code{copy_from_user} (při chybě \code{return -EFAULT})
    \item Nezapomeňte ukončit zkopírovaný řetězec použitím \code{\\0}
  \end{itemize}

  \item \code{read} bude vracet řetězec ,,Ahoj''
  \begin{itemize}
    \item \code{copy\_to\_user}
  \end{itemize}

  \end{enumerate}

\end{frame}

\begin{frame}
  \frametitle{Vytvoření znakového zařízení pomocí misc}

  \heading{Misc vrstva}
  \begin{itemize}
  \item \textbf{Stačí-li 1 minor} (málokdy)
  \item Dělá všechnu práci včetně volání \cmd{udev}

  \pause

  \item Potřebujeme
    \begin{itemize}
      \item Definici misc zařízení (\code{struct miscdevice}) -- minor číslo,
        jméno zařízení v \file{/dev}, atd.
      \item Opět \code{struct file_operations}
      \item \code{struct cdev} není třeba
    \end{itemize}
    \item \code{int misc_register(struct miscdevice *misc)}
    \item \code{void misc_deregister(struct miscdevice *misc)}
    \item Objeví se v \file{/proc/misc} a \file{/dev}
    \item \header{linux/miscdevice.h}
  \end{itemize}

\end{frame}

\begin{frame}
  \frametitle{Úkol}

  \heading{Vytvořit misc zařízení}
  \begin{enumerate}
    \item Definujte si nějakou \code{struct miscdevice}
    \begin{itemize}
      \item \code{minor = MISC_DYNAMIC_MINOR} (automaticky přidělit)
      \item \code{name} -- jméno vytvářeného zařízení
      \item \code{fops} -- vaše \code{struct file_operations}
    \end{itemize}

    \item \code{misc_register/deregister} do \code{module_init/exit}
    \item \cmd{make} a \cmd{insmod}
    \item Ověřte vytvoření v \file{/proc/misc}
    \item Vyzkoušejte \cmd{cat /dev/<name>} (jméno shora)
  \end{enumerate}
\end{frame}

\begin{frame}[fragile]
  \frametitle{Datové typy}
  \begin{itemize}
  \item Jádro a proces může běžet s různými bitovými šířkami (32, 64-bit)
  \item Problém s ukazateli a \code{long} proměnnými
  \begin{itemize}
  \item Jiná délka dat
  \item Jiné zarovnání struktur
  \end{itemize}
  \end{itemize}

  \bigskip

  \begin{center}
  \begin{tabular}{|p{2.8cm}|c|p{2.8cm}|} \hline
  Uživatel (32-bit) & & Jádro (64-bit) \\ \hline
  \code{data} \dots\ 4 bajty na offsetu 4 &
  \begin{lstlisting}[language=C]
struct my {
  char x;
  void *data;
};
  \end{lstlisting} &
  \code{data} \dots\ 8 bajtů na offsetu 8 \\ \hline
  \end{tabular}
  \end{center}

\end{frame}

\begin{frame}[fragile]
  \frametitle{Pevné typy v jádře}
  \framesubtitle{z hlediska počtu bitů}

  \begin{itemize}
  \item \header{linux/types.h}
  \item \code{__u8}, \code{__u16}, \code{__u32}, \code{__u64}
  \item \code{__s8}, \code{__s16}, \code{__s32}, \code{__s64}
  \item Ukazatele musí být v \code{union} s \code{__u64}
  \end{itemize}

  \begin{center}
  \begin{tabular}{ccc}
  \begin{lstlisting}
struct bad {
  unsigned long flags;
  short index;
  void *data;
} my;
  \end{lstlisting} & $\Rightarrow$ &
  \begin{lstlisting}
struct good {
  __u64 flags;
  __s16 index;
  union {
    void *data;
    __u64 filler;
  };
} my;
  \end{lstlisting}
  \end{tabular}
  \end{center}

  \begin{lstlisting}[columns=fullflexible]
read(fd, &my, sizeof(my));
ioctl(fd, DO_SOMETHING, &my);
  \end{lstlisting}

\end{frame}

